\newpage
\section{Current Computational Approaches and Limitations}
\label{sec2.2: current approaches}
% State of the Art in Uncertainty Quantification for Biological Prediction
The problem of predicting which phage infects which host bacteria has been studied computationally for over a 
decade and the methods have significantly evolved during this time. Prediction methods have progressed from simple rule 
based alignemnt to deep representation learning. Here, we will summarize each the primary methodologies of these methods
and along with the common limitations they share that motivate us for this work.


\subsection{Sequence Similarity and Protein Content Methods}
\label{sec2.2.1: classical methods}
The early host prediction tools rely on sequence alignment between a phage and bacterial genomes, assuming that 
similarity between the two indicates infection compatibility. HostPhinder \cite{villarroel2016hostphinder} represents each 
genome $g$ by its normalised $k$-mer frequency vector $\mathbf{v}(g) \in \mathbb{R}^{4^k}$ and assigns the predicted host 
through nearest-neighbour search over a reference database $\mathcal{D}_{\text{ref}} = \{(g_i, y_i)\}_{i=1}^N$:
\begin{equation*}
\hat{y} = \operatorname*{argmax}_{y \in \mathcal{H}} \max_{g_i:\, y_i = y} \operatorname{sim}\!\left(\mathbf{v}(g_{\text{test}}),\, \mathbf{v}(g_i)\right).
\end{equation*}

WIsH \cite{galiez2017wish} uses a probabilistic approach, it models each candidate host genome as a variable-
order Markov chain with parameters $\theta_h$ and then selects the host that maximises the log-likelihood of the new test 
phage sequence $S = (a_1, \ldots, a_L)$:
\begin{equation*}
\hat{y} = \operatorname*{argmax}_{h \in \mathcal{H}}\; \ell(S \mid \theta_h), \qquad \ell(S \mid \theta_h) = \sum_{i=k+1}
^{L} \log \mathbb{P}(a_i \mid a_{i-k}, \ldots, a_{i-1};\, \theta_h).
\end{equation*}

The protein-content methods overcome the limitation of whole-genome alignment by decomposing the phage into its 
constituent proteins. RaFAH \cite{coutinho2021rafah} clusters phage proteins into $M$ protein clusters via MMseqs2 
\cite{steinegger2017mmseqs2}, then it represents each phage by a binary membership vector $\mathbf{x} \in \{0,1\}^M$, 
and trains a random forest 
% $f_{\text{RF}}: \{0,1\}^M \to \Delta^{|\mathcal{H}|-1}$ 
to map this profile to a distribution over host genera. 
$$f_{\text{RF}}: \{0,1\}^M \longrightarrow \Delta^{|\mathcal{H}|-1}, \quad \text{where } \hat{\mathbf{p}} = 
    f_{\text{RF}}(\mathbf{x}), \quad \sum_{y \in \mathcal{H}} \hat{p}_y = 1$$
\noindent This representation is directly relevant to the present work, as our pipeline operates on a related 
protein-cluster construction, and the data contamination problem discussed in \ref{sec4.3: contamination fix} arises 
from the same source.

\subsection{Protein Language Models and Embedding-Based Methods}
\label{sec2.2.2: plms}
\noindent The current state of the art techniques use protein language models (pLMs), which are neural architectures 
trained on millions of protein sequences through self-supervised learning \cite{elnaggar2022prottrans, flamholz2024large}. 
Given a protein sequence $A = (a_1, \ldots, a_n) \in \mathcal{A}^n$ over the amino acid alphabet $\mathcal{A}$ of size 20, 
an encoder $\Phi_{\text{pLM}}$ produces a residue-level embedding matrix $\mathbf{Z} = \Phi_{\text{pLM}}(A) \in 
\mathbb{R}^{n \times d}$, it is then mean-pooled to a fixed-length sequence representation:
\begin{equation*}
\mathbf{z} = \frac{1}{n} \sum_{i=1}^{n} \mathbf{Z}_{i,:} \;\in\; \mathbb{R}^d.
\end{equation*}

The embedding $\mathbf{z}$ encodes biochemical, structural, and evolutionary information without requiring 
explicit sequence alignment. Building on ProtT5 \cite{elnaggar2022prottrans}, tools such as EMPATHI 
\cite{boulay2025empathi} and SUBLYME \cite{boulay2026phalp} use these representations for phage protein functional 
annotation and metagenomic lysin discovery respectively. These tools are directly involved in this work: the 
$K$-dimensional score vector $\mathbf{s} = \left( s_1(X), \dots, s_K(X) \right)^T\in [0,1]^K$ that serves as input to our 
conformal framework is produced by a pipeline built on these annotations, the precise construction of which is described 
in \ref{sec4: data}. For host prediction, protein embeddings have shown substantially improved accuracy over sequence-
similarity baselines \cite{gonzales2023protein, du2026high}. So this suggests that the embedding space contains sufficient 
information for reliable host assignment.


\subsection{Limitation: Absence of Uncertainty Quantification}
\label{sec2.2.3: limitations}
\noindent Despite their empirical success of all the methods seen, they all share a fundamental limitation: they produce 
point predictions (a single host label $\hat{y}$ or a scalar score $\hat{s}(x,y) \in [0,1]$) without any formal 
uncertainty quantification. 

\textbf{Uncertainty Quantification (UQ)} is the mathematical framework used to rigorously measure, 
propagate, and bound prediction errors. A raw scalar output $\hat{s}(x,y) = 0.87$ merely reflects an uncalibrated 
relative ranking statistic, not a true conditional probability $\mathbb{P}(Y = y \mid X = x)$. Existing host prediction 
tools fail to provide UQ because no tool offers a finite sample statistical guarantee that the true bacterial host is 
contained within a candidate prediction set with a predefined target probability $1-\alpha$.

Existing UQ frameworks only offer partial solutions. For example, the Bayesian approaches require strong assumptions 
on the distribution distributional that are difficult to justify for biological score data, while post-hoc calibration methods 
such as Platt scaling provide no finite-sample guarantees on unseen test data \cite{angelopoulos2023conformal}. 

These limitations are consequential in phage therapy, where an incorrectly predicted host can directly harm a 
patient. This motivates the adoption of conformal prediction \cite{vovk2005algorithmic}. This is a distribution-free 
framework that wraps any existing scoring model in a statistical construction, and produces prediction sets $C(X) 
\subseteq \mathcal{H}$ with exact finite-sample coverage guarantees while only requiring exchangeability of calibration 
and test data.
